\section{Louvain: Ordination, Preaching, and Controversial Method}

\subsection{A northern school in a divided Church}

Bellarmine was sent to Louvain in the spring of 1569.  The move situated a
young Italian Jesuit in a university and political world where Catholic and
Protestant claims were not remote textbook alternatives.  Teaching, worship,
preaching, civil order, and confessional allegiance intersected.  Bellarmine
continued his own study while being employed as a preacher and theological
teacher.  The surviving formation sequence does not warrant a claim that he
arrived as the holder of a completed modern theological degree.

He was ordained priest on 25 March 1570.  His autobiography remembers the
liturgical sequence leading through the orders to Holy Saturday; the present
civil date is supplied by official reception.  The exact place and complete
ordination dossier have not been collated for this study.  The result matters
more than ceremonial detail: preaching and theological exposition now became
priestly work ordered to sacramental and ecclesial service.

At Louvain Bellarmine cultivated Hebrew and worked through Scripture, patristic
authors, and the arguments of confessional opponents.  His
\emph{Institutiones linguae Hebraicae}, published in 1578 after his return to
Rome, witnesses that linguistic labor.  It does not make him a modern critical
philologist or guarantee the historical accuracy of every citation.  It does
show that his controversial method aspired to more than repeating scholastic
conclusions.  Text, history, inherited authority, and logical distinction had
to be assembled for public argument.

\subsection{Profession and a teaching identity}

Bellarmine professed the four Jesuit vows in July 1572.  Through 1576 he taught
and preached at Louvain while continuing to learn.  The dual role is essential
to his later work.  He did not first finish a private system and then descend
to controversy.  His theology developed in exchange with students,
congregations, printed opponents, and the institutional needs of a Church
receiving the Council of Trent amid continuing fracture.

Controversy in this setting was not automatically abuse.  It could identify a
real doctrinal difference, assemble the relevant authorities, distinguish
levels of claim, and show why a conclusion mattered for worship or communion.
It could also simplify opponents, intensify confessional boundaries, and make
victory seem equivalent to truth.  Bellarmine's mature work displays both the
strength of systematic confrontation and the moral risks of an institution
that possessed coercive as well as intellectual power.

The autobiography remembers Louvain as a formative success.  External
chronology supports the years and offices but cannot measure the private
development he assigned to them.  Later readers should therefore resist two
opposite reconstructions: that Bellarmine simply received an already complete
Roman program, or that he invented Catholic controversial theology alone.  He
worked within a Jesuit, university, patristic, biblical, and confessional
network whose demands gave his particular gifts a durable form.

\evidenceaxis{Bellarmine was ordained in 1570 and taught and preached at
Louvain until his recall to Rome in 1576.}{His autobiography controls the
broad sequence; modern critical and official records supply external dates;
his Hebrew grammar witnesses a continuing linguistic project.}{The exact
ordination place, matriculation, degree sequence, individual lectures, and
measure of his linguistic competence remain incompletely controlled.}
